
\subsection{Der singul\"are Ur-Fall $(5,7,11,13)$}

Der Vierling
\[
(5,7,11,13)
\]
nimmt innerhalb der Primzahlvierlinge eine Sonderstellung ein. Er ist
nicht nur der erste Vertreter der Form
\[
(p,p+2,p+6,p+8),
\]
sondern zugleich der einzige Fall, der nicht in das regul\"are
Vierlingsfenster
\[
(11,13,17,19)\pmod{30}
\]
f\"allt. Bereits statisch ist er daher kein blo\ss er Anfangspunkt einer
ansonsten homogenen Folge, sondern ein singul\"arer Grenzfall.

Diese Sonderstellung setzt sich dynamisch fort. Unter der reduzierten
Collatz-Abbildung
\[
T(n)=\frac{3n+1}{2^{v_2(3n+1)}}
\]
ergibt sich
\[
(5,7,11,13)\mapsto(1,11,17,5),
\]
mit Kompressionsprofil
\[
V=(4,1,1,3).
\]
Schon im ersten Schritt fallen damit alle vier Komponenten in einen sehr
kleinen ungeraden Kernraum. Auch die folgenden Bilder bleiben extrem
kompakt:
\[
(1,11,17,5)\mapsto(1,17,13,1)\mapsto(1,13,5,1)\mapsto(1,5,1,1).
\]
Die fr\"uhe Spreizung bleibt klein und sinkt rasch ab; zugleich ist der
R\"uckkonzentrationsindex bereits vom ersten Schritt an maximal.

Modellhaft kann dieser Sachverhalt als besonders dicht gebundenes
Ur-Tetraeder gelesen werden. W\"ahrend die sp\"ateren Primzahlvierlinge
zwar ebenfalls saturierte lokale Schlie\ss ungszellen darstellen, aber
unter $T$ in vier deutlich getrennte Kernlinien entb\"undelt werden,
bleibt der Anfangsfall von Anfang an in unmittelbarer N\"ahe des
minimalen Kerns. Gerade deshalb besitzt er innerhalb der Theorie nicht
nur historischen, sondern strukturellen Rang: Er ist nicht einfach der
erste Vierling, sondern der kompakteste.
